\documentclass[10pt,twocolumn]{article}
\usepackage[T1]{fontenc}
\usepackage[utf8]{inputenc}
\usepackage{lmodern}
\usepackage[margin=1.8cm,top=2.4cm,bottom=2cm,columnsep=0.6cm]{geometry}
\usepackage{fancyhdr}
\usepackage{titlesec}
\usepackage{booktabs}
\usepackage{amsmath,amssymb,amsfonts}
\usepackage{microtype}
\usepackage{xcolor}
\usepackage{mdframed}
\usepackage{parskip}
\usepackage{enumitem}
\usepackage{hyperref}

\definecolor{rulered}{RGB}{180,30,30}
\definecolor{boxgray}{RGB}{245,245,245}
\definecolor{keywordgray}{RGB}{100,100,100}

\pagestyle{fancy}
\fancyhf{}
\fancyhead[L]{\textsc{\small Claude Mag}}
\fancyhead[C]{\small\textit{Machine Learning Research Digest}}
\fancyhead[R]{\small 2026-10-07}
\fancyfoot[C]{\small\thepage}
\renewcommand{\headrulewidth}{0.8pt}

\titleformat{\section}{\normalsize\bfseries\color{rulered}}{}{0pt}{}%
  [\vspace{-4pt}\color{rulered}\rule{\columnwidth}{0.4pt}]
\titlespacing{\section}{0pt}{8pt}{4pt}

\hypersetup{colorlinks=true,urlcolor=rulered,linkcolor=black}

\begin{document}
\setlength{\parindent}{0pt}
\setlength{\parskip}{4pt}

{\small\color{keywordgray}\textbf{Keywords:}
  token cues \textbullet{}
  base models \textbullet{}
  reasoning \textbullet{}
  causal data intervention}\par\vspace{4pt}

{\LARGE\bfseries\raggedright
  Base Models Can Reason By Taking a Cue From Training Data}\par\vspace{4pt}

{\large\itshape\raggedright
  A Single Starting Token Can Raise Math Accuracy by 36~pp in a Base Model---and
  the Effect Is Causally Traceable to Training Data, Not Model Architecture}\par\vspace{6pt}

{\small\textbf{By} Sophie L.\ Wang, Amil Dravid, Rulin Shao, Kevin Farhat,
  Sewon Min \& Alexei A.\ Efros (UC Berkeley)}\par\vspace{2pt}

{\small\color{keywordgray}arXiv:2610.06851 \textbullet{} October 2026}
\par\vspace{2pt}\noindent{\color{rulered}\rule{\columnwidth}{0.6pt}}\vspace{6pt}

Base language models harbor latent reasoning competence that is largely dormant until
activated by a specific starting token---a \emph{reasoning cue} whose effect is
not semantic but distributional, mechanically traceable to the cluster of
step-by-step reasoning documents in the training corpus that bear the same token.

\begin{mdframed}[backgroundcolor=boxgray,linecolor=rulered,linewidth=1pt,
    innerleftmargin=6pt,innerrightmargin=6pt,
    innertopmargin=6pt,innerbottommargin=6pt]
\textbf{\small AT A GLANCE}\par\vspace{4pt}
\begin{description}[leftmargin=0pt,labelsep=4pt,itemsep=2pt,parsep=0pt,
    font=\small\bfseries]
  \item[Problem:] Reinforcement learning post-training is costly; the mechanism
    by which it elicits reasoning chains is poorly understood.
  \item[Why it matters:] If RL merely selects regions of a base model's training
    distribution, cheaper interventions should reproduce the gains.
  \item[Method:] Systematic sweep over starting-token cues; causal data
    interventions to add or remove reasoning associations.
  \item[Result:] The cue ``.\textbackslash n\textbackslash nOkay'' raises Olmo-3-7B MATH-500 from 42\%
    to 78\%; causal interventions confirm the training-data mechanism.
  \item[Evidence:] Evaluated on MATH-500, coding benchmarks, and multiple
    frontier base models; causal interventions with data ablations.
\end{description}
\end{mdframed}

\section{The Big Picture}

Large language models trained with reinforcement learning from verifiable rewards
(RLVR) achieve dramatic performance gains on mathematical reasoning and code
generation. Yet the precise mechanism has remained unclear: does RL teach the
model new reasoning skills, or does it teach the model \emph{when} to deploy
skills that the base model already possesses?

This paper proposes a third, more concrete answer: RL teaches the model to begin
its response with particular starting tokens---reasoning cues---that route
generation toward the region of the training distribution in which careful
step-by-step reasoning documents reside. The paper provides both correlational
and causal evidence for this mechanism.

\section{Why Existing Approaches Fall Short}

\textbf{RL post-training is expensive.} RLVR requires generating many rollouts and
computing rewards, multiplying training compute relative to supervised fine-tuning.
If the mechanism is a distributional routing effect, a cheaper intervention should
suffice.

\textbf{Prompt engineering is empirical.} Prior chain-of-thought work shows that
phrases like ``Let's think step by step'' improve performance, but attributes
the effect to semantic instruction following---a hypothesis that predicts the
effect should generalize to semantically equivalent phrases, which the paper tests
and finds false.

\textbf{Base model latent capacity is underutilized.} The default starting token
chosen by a practitioner is arbitrary; this work shows it has outsized influence
on the quality of the continuation.

\section{Core Mechanism}

A \textbf{reasoning cue} is a short prefix $c$ placed at the start of a base model's
response. The central finding is that effective cues are not semantically special:
their power comes from co-occurring in the training corpus with step-by-step
reasoning chains, not from the meaning of the words themselves.

When a cue $c$ appears at the start of the model's response, next-token prediction
conditions on $c$ and selects continuations from the region of the training data
distribution where $c$ was associated with careful, multi-step reasoning.
This makes the cue a \emph{distributional selector} rather than a semantic instruction.

\section{Technical Formulation}

Let $\mathcal{M}$ be a base LLM with vocabulary $\mathcal{V}$. For a problem $q$
and cue string $c \in \mathcal{V}^*$:
\[
\text{Acc}(c) = \mathbb{E}_q\bigl[\text{Correct}(\mathcal{M}(q, c))\bigr]
\]
The sweep identifies the optimal cue $c^* = \arg\max_c \text{Acc}(c)$.

For the causal intervention, let $D_c \subseteq \mathcal{D}_\text{train}$ be
training documents whose response begins with $c$. The paper shows:
\[
\text{Acc}_{D \setminus D_c}(c) \ll \text{Acc}_D(c)
\]
Conversely, augmenting the training data with arbitrary-token-prefixed reasoning
traces $\widetilde{D}_{c'}$ creates a new effective cue:
\[
\text{Acc}_{D \cup \widetilde{D}_{c'}}(c') \approx \text{Acc}_D(c^*)
\]

\section{Learning or Inference Procedure}

\textbf{Finding effective cues:}
\begin{enumerate}[itemsep=2pt,parsep=0pt]
  \item Enumerate a large set of candidate starting tokens and phrases.
  \item For each cue, generate responses on a held-out problem set.
  \item Score responses and rank cues by mean accuracy.
\end{enumerate}

\textbf{Causal confirmation:}
\begin{enumerate}[itemsep=2pt,parsep=0pt]
  \item Identify training documents co-occurring with the effective cue.
  \item Remove those documents and retrain (or fine-tune) a model.
  \item Verify that the cue's effect disappears with the associated data.
\end{enumerate}

\textbf{Creating custom cues:}
\begin{enumerate}[itemsep=2pt,parsep=0pt]
  \item Collect reasoning traces; prepend an arbitrary token to each.
  \item Fine-tune on these modified traces.
  \item The arbitrary token becomes an effective cue.
\end{enumerate}

\section{What the Guarantee Says}

The guarantee is empirical: across all tested base models and cue candidates,
the reasoning benefit of any cue is causally traceable to its associated training
documents. Removing those documents degrades the cue; adding documents that
associate an arbitrary token with reasoning creates a new effective cue.
The mechanism is not model-specific---it holds for Olmo-3-7B, Qwen3-14B, and
several other frontier base models.

\section{Experimental Findings}

\begin{itemize}[itemsep=2pt,parsep=0pt]
  \item \textbf{Olmo-3-7B:} the cue ``.\textbackslash n\textbackslash nOkay'' raises
    MATH-500 pass@1 from \textbf{42\% to 78\%}, matching the RL-trained counterpart.
  \item \textbf{Qwen3-14B:} ``Alright,'' raises MATH-500 from \textbf{72\% to 87\%},
    also competitive with RL.
  \item \textbf{Code generation:} analogous cues improve pass@$k$ metrics on
    coding benchmarks.
  \item \textbf{Causal intervention:} removing documents associated with the
    effective cue eliminates its accuracy benefit.
  \item \textbf{Custom cue creation:} the word ``chicken,'' after fine-tuning on
    reasoning traces prefixed with it, becomes an effective cue.
\end{itemize}

\section{Ablations and Interpretation}

\textbf{Semantic vs.\ distributional hypothesis:} Replacing the effective cue with
semantically similar phrases that do not appear in the relevant training data
produces no reasoning gain, confirming the distributional mechanism.

\textbf{Context length stability:} The cue effect holds even with few-shot examples
in context, ruling out attention-dilution as an alternative explanation.

\textbf{Multiple cues:} Combining multiple effective cues does not additively
compound gains, indicating they tap overlapping training data regions.

\textbf{Implications for RL:} Since base-model cue performance matches RLVR, the
paper argues that RL primarily teaches models to start responses in
the step-by-step reasoning region of their training distribution---not to
acquire fundamentally new reasoning capabilities.

\section*{Reference}

Sophie L.\ Wang et al.\
\textbf{Base Models Can Reason By Taking a Cue From Training Data}.
arXiv:2610.06851, October 2026.\\
\url{https://arxiv.org/abs/2610.06851}

\end{document}
